% All drawings are native, reproducible TikZ vectors.  Geometry panels use
% identical x/y units; panels marked schematic depict decisions, not a claimed
% complete subdivision of a particular input.
\tikzset{coordgrid/.style={draw=quiet!24,line width=.25pt}}

% Reading map: recursive actions and their queries, not a claimed cell proof.
\newcommand{\figoverview}{%
\begin{tikzpicture}[font=\small,
 overview/.style={draw=navy!40,fill=codebg,rounded corners=2pt,
   text width=12.5cm,align=center,inner sep=8pt},
 action/.style={draw=navy!40,fill=white,rounded corners=2pt,
   text width=3.65cm,minimum height=3.3cm,align=center,inner sep=7pt}]
\node[overview] (prep) at (0,0) {%
  {\bfseries\color{navy}Prepare the level-\(i\) boundary}\\[2pt]
  Normalize \(P_i\) and split its edges at earlier boundaries and at \(s\)
  when \(s\) lies on an edge.
  Build needed incident rays from two one-sided queries at level \(i-1\)
  (lazily or eagerly).};
\node[overview,fill=good!6,draw=good!55] (loc) at (0,-3.1) {%
  {\bfseries\color{good}At \((q,i)\): choose the last-step action}\\[2pt]
  Test closed membership first. For an exterior query, locate a vertex cone
  or following pseudo-edge using the rays and chord predicates.};
\node[action,fill=good!6] (crossing) at (-5,-7.4) {%
  {\bfseries\color{good}CROSS}\\[5pt]
  The previous optimal route already supplies visit \(i\).\\[7pt]
  \(F_i(q)=F_{i-1}(q)\)\\[7pt]
  Recurse at \(q\).};
\node[action,fill=pOne!6] (bend) at (0,-7.4) {%
  {\bfseries\color{pOne}BEND AT \(v\)}\\[5pt]
  The last chosen contact is a split vertex.\\[7pt]
  \(F_i(q)=F_{i-1}(v)+\|q-v\|\)\\[7pt]
  Recurse at \(v\); append \(vq\).};
\node[action,fill=pTwo!7] (reflect) at (5,-7.4) {%
  {\bfseries\color{pTwo}REFLECT IN \(e\)}\\[5pt]
  Unfold through the parent edge line.\\[7pt]
  \(q'=R_e(q),\quad F_i(q)=F_{i-1}(q')\)\\[7pt]
  Recurse at \(q'\); refold on finite \(e\).};
\node[overview] (base) at (0,-11.8) {%
  {\bfseries\color{navy}Return from the recursion}\\[2pt]
  At level \(0\), \(F_0(q)=\|q-s\|\). Reconstruct contacts in return order;
  a path query checks the finite-edge intersections before compacting
  forward-collinear points.};
\draw[quiet,-Latex,line width=.9pt] (prep.south)--(loc.north);
\foreach \target in {crossing,bend,reflect}
  \draw[quiet,-Latex,line width=.9pt] (loc.south)--(\target.north);
\foreach \target in {crossing,bend,reflect}
  \draw[quiet,-Latex,line width=.9pt] (\target.south)--(base.north);
\end{tikzpicture}}

\newcommand{\figreflection}{%
\begin{tikzpicture}[scale=.92]
\begin{scope}
  \draw[coordgrid] (-3,-1) grid[step=.5] (3,3);
  \draw[polyA] (0,0) rectangle (2,2);
  \draw[very thick,navy] (0,0)--(0,2);
  \draw[route,-Latex] (-2,1)--(0,1/3);
  \draw[route,-Latex] (0,1/3)--(-1,0);
  \fill[ink] (-2,1) circle (1.8pt);
  \fill[ink] (-1,0) circle (1.8pt);
  \fill[good] (0,1/3) circle (2.2pt);
  \node[smalllabel,above left] at (-2,1) {\(s\)};
  \node[smalllabel,below] at (-1,-.08) {\(q\)};
  \node[smalllabel,right] at (.08,.38) {\(c\)};
  \node[smalllabel] at (1.1,1.1) {\(P\)};
  \node[font=\scriptsize\bfseries,color=navy] at (0,-1.52)
    {(a) physical route};
\end{scope}
\begin{scope}[xshift=7.4cm]
  \draw[coordgrid] (-3,-1) grid[step=.5] (3,3);
  \draw[very thick,navy] (0,0)--(0,2);
  \draw[route,-Latex] (-2,1)--(1,0);
  \draw[wrong,dashed] (-1,0)--(1,0);
  \fill[ink] (-2,1) circle (1.8pt);
  \fill[bad] (-1,0) circle (1.8pt);
  \fill[good] (0,1/3) circle (2.2pt);
  \fill[ink] (1,0) circle (1.8pt);
  \node[smalllabel,above left] at (-2,1) {\(s\)};
  \node[smalllabel,below] at (-1,-.08) {\(q\)};
  \node[smalllabel,right] at (.08,.38) {\(c\)};
  \node[smalllabel,below] at (1,-.08) {\(q'\)};
  \node[font=\scriptsize\bfseries,color=navy] at (0,-1.52)
    {(b) unfolded segment};
\end{scope}
\end{tikzpicture}}

\newcommand{\figcrossbend}{%
\begin{tikzpicture}[scale=1.12]
\begin{scope}
  \draw[polyA] (0,0) rectangle (2,2);
  \node[smalllabel] at (1,1) {\(P\)};
  \draw[coordgrid] (-2,-2) grid[step=.5] (3,2);
  \draw[route,-Latex] (-2,1)--(3,1);
  \fill[ink] (-2,1) circle (1.8pt); \fill[ink] (3,1) circle (1.8pt);
  \node[smalllabel,above left] at (-2,1) {\(s\)};
  \node[smalllabel,above right] at (3,1) {\(q\)};
  \node[font=\scriptsize\bfseries,color=navy,align=center] at (1,2.65)
    {crossing\\\(s=(-2,1),\ q=(3,1)\)};
  \node[smalllabel,align=center] at (1,-2.58)
    {straight route intersects \(P\)};
\end{scope}
\begin{scope}[xshift=5.5cm]
  \draw[polyA] (0,0) rectangle (2,2);
  \node[smalllabel] at (1,1) {\(P\)};
  \draw[coordgrid] (-2,-2) grid[step=.5] (2,2);
  \draw[route,-Latex] (-2,-1)--(0,0);
  \draw[route,-Latex] (0,0)--(1,-2);
  \fill[ink] (-2,-1) circle (1.8pt); \fill[ink] (1,-2) circle (1.8pt);
  \fill[good] (0,0) circle (2.2pt);
  \node[smalllabel,below left] at (-2,-1) {\(s\)};
  \node[smalllabel,below right] at (1,-2) {\(q\)};
  \node[smalllabel,above right] at (0,0) {\(v=(0,0)\)};
  \node[font=\scriptsize\bfseries,color=navy,align=center] at (1,2.65)
    {vertex bend\\\(s=(-2,-1),\ q=(1,-2)\)};
  \node[smalllabel] at (1,-2.58) {length \(2\sqrt5\)};
\end{scope}
\end{tikzpicture}}

\newcommand{\figrepeat}{%
\begin{tikzpicture}[scale=.78]
\begin{scope}
\draw[polyA] (2,1) rectangle (4,3);
\node[smalllabel] at (3,2.45) {\(P_1=P_2=P_3\)};
\draw[route,-Latex] (0,2)--(6,2);
\fill[ink] (0,2) circle (1.8pt); \fill[ink] (6,2) circle (1.8pt);
\fill[good] (2,2) circle (2.2pt);
\node[smalllabel,above] at (0,2.12) {\(s\)};
\node[smalllabel,above] at (6,2.12) {\(q\)};
\node[smalllabel,below] at (2,1.9) {visits \(1=2=3\)};
\node[font=\scriptsize\bfseries,color=navy] at (3,-.55)
  {(a) \(sq\) passes through the three squares};
\end{scope}
\begin{scope}[xshift=7.3cm]
\draw[polyA] (2,1) rectangle (4,3);
\node[smalllabel] at (3,2.45) {\(P_1=P_2=P_3\)};
\draw[route,-Latex] (0,2)--(3,2);
\fill[ink] (0,2) circle (1.8pt); \fill[good] (3,2) circle (2.2pt);
\node[smalllabel,above] at (0,2.12) {\(s\)};
\node[smalllabel,below] at (3,1.9) {\(q:\ 1=2=3\)};
\node[font=\scriptsize\bfseries,color=navy] at (3,-.55)
  {(b) \(q\) lies in all three squares};
\end{scope}
\end{tikzpicture}}

\newcommand{\figfalsecrossing}{%
\begin{tikzpicture}[scale=.63]
\begin{scope}
  \draw[polyA] (0,2) rectangle (2,4);
  \draw[polyB] (0,-2) rectangle (3,2);
  \draw[draw=good,fill=good!10,line width=.9pt] (-1,3) rectangle (1,4);
  \draw[coordgrid] (-4,-4) grid[step=1] (4,4);
  \draw[navy!70,-Latex,line width=.9pt] (-4,0)--(4.35,0)
    node[smalllabel,right] {\(x\)};
  \draw[navy!70,-Latex,line width=.9pt] (0,-4)--(0,4.35)
    node[smalllabel,above] {\(y\)};
  \node[smalllabel] at (1.65,2.7) {\(P_1\)};
  \node[smalllabel] at (2.25,.35) {\(P_2\)};
  \node[smalllabel] at (-.65,3.7) {\(P_3\)};
  \fill[ink] (-4,-4) circle (2pt); \fill[ink] (4,-4) circle (2pt);
  \node[smalllabel,below] at (-4,-4) {\(s\)};
  \node[smalllabel,below] at (4,-4) {\(t\)};
  \draw[wrong,-Latex] (-4,-4)--(0,3);
  \draw[wrong,-Latex] (0,3)--(4,-4);
  \fill[bad] (0,3) circle (2.4pt);
  \fill[bad] (4/7,2) circle (2.4pt);
  \node[smalllabel,left] at (0,3) {\(1,3\)};
  \node[smalllabel,right] at (4/7,2) {\(2\) later};
  \node[smalllabel,align=center] at (0,-5.35)
    {invalid: \(P_1\) and \(P_3\) together, then \(P_2\)};
  \node[font=\footnotesize\bfseries,color=navy] at (0,5.15) {(a) old route};
\end{scope}
\begin{scope}[xshift=11cm]
  \draw[polyA] (0,2) rectangle (2,4);
  \draw[polyB] (0,-2) rectangle (3,2);
  \draw[draw=good,fill=good!10,line width=.9pt] (-1,3) rectangle (1,4);
  \draw[coordgrid] (-4,-4) grid[step=1] (4,4);
  \draw[navy!70,-Latex,line width=.9pt] (-4,0)--(4.35,0)
    node[smalllabel,right] {\(x\)};
  \draw[navy!70,-Latex,line width=.9pt] (0,-4)--(0,4.35)
    node[smalllabel,above] {\(y\)};
  \node[smalllabel] at (1.65,2.7) {\(P_1\)};
  \node[smalllabel] at (2.25,.35) {\(P_2\)};
  \node[smalllabel] at (-.65,3.7) {\(P_3\)};
  \fill[ink] (-4,-4) circle (2pt); \fill[ink] (4,-4) circle (2pt);
  \node[smalllabel,below] at (-4,-4) {\(s\)};
  \node[smalllabel,below] at (4,-4) {\(t\)};
  \draw[route,-Latex] (-4,-4)--(0,2);
  \draw[route,-Latex] (0,2)--(.5,3);
  \draw[route,-Latex] (.5,3)--(4,-4);
  \fill[good] (0,2) circle (2.4pt); \fill[good] (.5,3) circle (2.4pt);
  \node[smalllabel,left] at (0,2) {\(1=2\)};
  \node[smalllabel,right] at (.5,3) {\(3\)};
  \node[smalllabel,align=center] at (0,-5.35)
    {valid: \(P_1\) and \(P_2\) together, then \(P_3\)};
  \node[font=\footnotesize\bfseries,color=navy] at (0,5.15) {(b) corrected route};
\end{scope}
\end{tikzpicture}}

\newcommand{\figfalsecrossingzoom}{%
\begin{tikzpicture}[scale=2.65]
\begin{scope}
  \clip (-1.15,1.35) rectangle (1.65,3.65);
  \draw[coordgrid] (-2,1) grid[step=.5] (2,4);
  \draw[polyB] (0,-2) rectangle (3,2);
  \draw[polyA] (0,2) rectangle (2,4);
  \draw[draw=good,fill=good!10,line width=.9pt] (-1,3) rectangle (1,4);
  \draw[wrong,-Latex,line width=1.1pt] (-4,-4)--(0,3);
  \draw[wrong,-Latex,line width=1.1pt] (0,3)--(4,-4);
  \draw[route,-Latex] (-4,-4)--(0,2);
  \draw[route,-Latex] (0,2)--(.5,3);
  \draw[route,-Latex] (.5,3)--(4,-4);
\end{scope}
\draw[navy!45] (-1.15,1.35) rectangle (1.65,3.65);
\fill[bad] (0,3) circle (1.5pt); \fill[bad] (4/7,2) circle (1.5pt);
\fill[good] (0,2) circle (1.5pt); \fill[good] (.5,3) circle (1.5pt);
\node[smalllabel,above left] at (0,3) {old: \(1,3\)};
\node[smalllabel,right] at (4/7,2) {old: \(2\)};
\node[smalllabel,left] at (0,2) {valid: \(1,2\)};
\node[smalllabel,right] at (.5,3) {valid: \(3\)};
\node[smalllabel] at (-.75,3.4) {\(P_3\)};
\node[smalllabel] at (1.3,3.4) {\(P_1\)};
\node[smalllabel] at (1.3,1.65) {\(P_2\)};
\end{tikzpicture}}

\newcommand{\figordered}{%
\begin{tikzpicture}[scale=.72]
\begin{scope}
  \path (-.5,-.45) rectangle (8.5,3);
  \draw[polyB] (2,1.5) rectangle (3,2.5);
  \draw[polyA] (5,1.5) rectangle (6,2.5);
  \node[smalllabel] at (2.5,2.05) {\(P_2\)};
  \node[smalllabel] at (5.5,2.05) {\(P_1\)};
  \draw[route,-Latex] (0,0)--(5,1.5);
  \draw[route,-Latex] (5,1.5)--(3,1.5);
  \draw[route,-Latex] (3,1.5)--(8,0);
  \foreach \x/\y/\lab in {0/0/s,8/0/t}
    {\fill[ink] (\x,\y) circle (2pt);
     \node[smalllabel,below] at (\x,\y-.08) {\(\lab\)};}
  \node[font=\scriptsize\bfseries,color=good] at (4,.42) {feasible};
\end{scope}
\begin{scope}[xshift=10.7cm]
  \path (-.5,-.45) rectangle (8.5,3);
  \draw[polyB] (2,1.5) rectangle (3,2.5);
  \draw[polyA] (5,1.5) rectangle (6,2.5);
  \node[smalllabel] at (2.5,2.05) {\(P_2\)};
  \node[smalllabel] at (5.5,2.05) {\(P_1\)};
  \draw[wrong,-Latex,line width=1.2pt] (0,0)--(3,1.5);
  \draw[wrong,-Latex,line width=1.2pt] (3,1.5)--(5,1.5);
  \draw[wrong,-Latex,line width=1.2pt] (5,1.5)--(8,0);
  \foreach \x/\y/\lab in {0/0/s,8/0/t}
    {\fill[ink] (\x,\y) circle (2pt);
     \node[smalllabel,below] at (\x,\y-.08) {\(\lab\)};}
  \node[font=\scriptsize\bfseries,color=bad] at (4,.42) {infeasible};
\end{scope}
\end{tikzpicture}}
\newcommand{\figcontacts}{%
\begin{tikzpicture}[scale=.82]
\begin{scope}
  \draw[polyA] (0,0) rectangle (2,2);
  \node[smalllabel] at (1,1.1) {\(P_1\)};
  \draw[route,-Latex] (-2,-1)--(3,2);
  \fill[ink] (-2,-1) circle (2pt); \fill[ink] (3,2) circle (2pt);
  \node[smalllabel,below left] at (-2,-1) {\(s\)};
  \node[smalllabel,above right] at (3,2) {\(q_C\)};
  \node[font=\small\bfseries,color=navy] at (.5,3) {crossing};
\end{scope}
\begin{scope}[xshift=6.4cm]
  \draw[polyA] (0,0) rectangle (2,2);
  \node[smalllabel] at (1,1.1) {\(P_1\)};
  \draw[route,-Latex] (-2,-1)--(0,0);
  \draw[route,-Latex] (0,0)--(1,-2);
  \fill[ink] (-2,-1) circle (2pt); \fill[ink] (1,-2) circle (2pt);
  \fill[good] (0,0) circle (2.2pt);
  \node[smalllabel,below left] at (-2,-1) {\(s\)};
  \node[smalllabel,below right] at (1,-2) {\(q_V\)};
  \node[smalllabel,above left] at (0,0) {\(v\)};
  \node[font=\small\bfseries,color=navy] at (.5,3) {vertex bend};
\end{scope}
\begin{scope}[xshift=12.8cm]
  \draw[polyA] (0,0) rectangle (2,2);
  \node[smalllabel] at (1,1.1) {\(P_1\)};
  \draw[route,-Latex] (-2,-1)--(.5,0);
  \draw[route,-Latex] (.5,0)--(3,-1);
  \fill[ink] (-2,-1) circle (2pt); \fill[ink] (3,-1) circle (2pt);
  \fill[good] (.5,0) circle (2.2pt);
  \node[smalllabel,below left] at (-2,-1) {\(s\)};
  \node[smalllabel,below right] at (3,-1) {\(q_E\)};
  \node[smalllabel,above] at (.5,0) {\(c\)};
  \node[font=\small\bfseries,color=navy] at (.5,3) {edge reflection};
\end{scope}
\end{tikzpicture}}
\newcommand{\figunfold}{%
\begin{tikzpicture}[scale=1.08]
\draw[polyA] (-.3,0)--(4.2,0); \node[smalllabel] at (4, .45) {\(e\)};
\coordinate (S) at (.2,-2); \coordinate (C) at (2,0); \coordinate (Q) at (3.8,-2); \coordinate (R) at (3.8,2);
\draw[route,-Latex] (S)--(C)--(Q); \draw[aux,-Latex] (S)--(R);
\draw[aux] (Q)--(R); \fill[ink] (C) circle (1.5pt);
\node[smalllabel,below] at (S) {\(a\)}; \node[smalllabel,below] at (Q) {\(q\)};
\node[smalllabel,above] at (R) {\(q'=R_e(q)\)};
\node[smalllabel,above left] at (C) {\(c\)};
\end{tikzpicture}}
\newcommand{\figlastmap}{%
\begin{tikzpicture}[scale=1.24]
\begin{scope}
  \clip (-4,-2.4) rectangle (4,4.4);
  \fill[bad!17] (-4,-2.4)--(4,-2.4)--(4,-2)--(0,0)--(-4,-2)--cycle;
  \fill[bad!17] (-4,4.4)--(4,4.4)--(4,4)--(0,2)--(-4,4)--cycle;
  \fill[pTwo!22] (-4,-2)--(0,0)--(0,2)--(-4,4)--cycle;
  \fill[good!16] (0,0)--(4,-2)--(4,4)--(0,2)--cycle;
  \draw[polyA,fill=pOne!27] (0,0) rectangle (2,2);
  \draw[aux,line width=.9pt] (-4,-2)--(0,0)--(4,-2);
  \draw[aux,line width=.9pt] (-4,4)--(0,2)--(4,4);
  \draw[very thick,navy] (0,0)--(0,2);
\end{scope}
\draw[navy!45] (-4,-2.4) rectangle (4,4.4);
\fill[ink] (-2,1) circle (2.3pt);
\node[smalllabel,above] at (-2,1) {\(s=(-2,1)\)};
\node[smalllabel] at (1,1.25) {\(P_1\)};
\node[smalllabel,font=\footnotesize\bfseries,text=ink] at (0,3.15) {bend at \((0,2)\)};
\node[smalllabel,font=\footnotesize\bfseries,text=ink] at (0,-1.3) {bend at \((0,0)\)};
\node[smalllabel,font=\footnotesize\bfseries,text=ink] at (-2,-.45) {left-edge reflection};
\node[smalllabel,font=\footnotesize\bfseries,text=ink] at (2.5,.4) {cross};
\node[smalllabel,below right] at (0,0) {\((0,0)\)};
\node[smalllabel,above right] at (0,2) {\((0,2)\)};
\draw[good,-Latex,line width=.8pt] (-2,1)--(3,1);
\draw[pTwo,-Latex,line width=.8pt] (-2,1)--(0,.6);
\draw[pTwo,-Latex,line width=.8pt] (0,.6)--(-3,0);
\draw[bad,-Latex,line width=.8pt] (-2,1)--(0,2);
\draw[bad,-Latex,line width=.8pt] (0,2)--(1,4);
\draw[bad,-Latex,line width=.8pt] (-2,1)--(0,0);
\draw[bad,-Latex,line width=.8pt] (0,0)--(1,-2);
\foreach \x/\y/\lab in {3/1/C,-3/0/R,1/4/U,1/-2/L}
  {\fill[ink] (\x,\y) circle (2.2pt);
   \node[smalllabel,right] at (\x+.08,\y) {\(q_{\lab}\)};}
\end{tikzpicture}}

\newcommand{\figwitnesssetup}{%
\begin{tikzpicture}[scale=.64]
  \path[fill=pTwo!12] (-8,-4)--(8,4)--(8,9)--(-8,9)--cycle;
  \path[fill=pOne!12] (0,-6) rectangle (6,6);
  \draw[coordgrid] (-9,-7) grid[step=2] (9,10);
  \draw[navy!70,-Latex,line width=1.2pt] (-9,0)--(9.35,0)
    node[smalllabel,right] {\(x\)};
  \draw[navy!70,-Latex,line width=1.2pt] (0,-7)--(0,10.35)
    node[smalllabel,above] {\(y\)};
  \draw[polyB,fill=none] (-8,-4)--(8,4)--(8,9)--(-8,9)--cycle;
  \draw[polyA,fill=none] (0,-6) rectangle (6,6);
  \draw[very thick,navy] (0,-6)--(0,6);
  \draw[very thick,pTwo] (-8,-4)--(8,4);
  \node[smalllabel] at (4.8,5.2) {\(P_1\)};
  \node[smalllabel] at (-5.8,6.5) {\(P_2\)};
  \foreach \x/\y in {-3/-4,0/0,4/-3,-4/-3,6/3}
    {\fill[ink] (\x,\y) circle (2.3pt);}
  \foreach \x/\y in {0/-6,6/6,-8/-4,8/4}
    {\fill[ink] (\x,\y) circle (2.3pt);}
  \node[smalllabel,below right] at (-3,-4) {\(s=(-3,-4)\)};
  \node[smalllabel,above left] at (-4,-3) {\(t_B=(-4,-3)\)};
  \node[smalllabel,below right] at (4,-3) {\(t_A=(4,-3)\)};
  \node[smalllabel,above left] at (0,0) {\(j=(0,0)\)};
  \node[smalllabel,above right] at (6,3) {\((6,3)\)};
  \node[smalllabel,below right] at (0,-6) {\((0,-6)\)};
  \node[smalllabel,above left] at (6,6) {\((6,6)\)};
  \node[smalllabel,below right] at (-8,-4) {\((-8,-4)\)};
  \node[smalllabel,above left] at (8,4) {\((8,4)\)};
  \node[smalllabel,above] at (-1,9) {\(y=9\)};
  \node[smalllabel,right] at (0,-5.2) {\(\ell_1:x=0\)};
  \node[smalllabel,below right] at (4,2) {\(\ell_2:y=x/2\)};
\end{tikzpicture}}
\newcommand{\figbinary}{%
\begin{tikzpicture}
\foreach \shift/\j/\endx/\endy/\rayx/\rayy/\txt in
 {0/3/8/4/2.6/.8/{keep \(e_0,e_1,e_2\)},
  5.3/2/6/3/2.2/.6/{keep \(e_0,e_1\)},
  10.6/1/0/0/.7/-2.4/{keep \(e_0\)}} {
  \begin{scope}[xshift=\shift cm,scale=.27]
    \draw[polyB] (-8,-4)--(0,0)--(6,3)--(8,4)--(8,9)--(-8,9)--cycle;
    \node[smalllabel] at (-4,5) {\(P_2\)};
    \draw[bad,line width=2pt] (\endx,\endy)--(8,4)--(8,9)--(-8,9)--(-8,-4);
    \draw[good,line width=2pt] (-8,-4)--(\endx,\endy);
    \draw[aux,line width=1.2pt] (-8,-4)--(\endx,\endy);
    \draw[good,-Latex] (\endx,\endy)--++(\rayx,\rayy);
    \fill[ink] (-8,-4) circle (4pt);
    \fill[ink] (\endx,\endy) circle (4pt);
    \fill[good] (-4,-3) circle (5pt);
    \node[font=\scriptsize,below left] at (-8,-4) {\(v_0\)};
    \node[font=\scriptsize,above right] at (\endx,\endy) {\(v_{\j}\)};
    \node[font=\scriptsize,below] at (-4,-3) {\(q\)};
  \end{scope}
  \node[font=\footnotesize,align=center] at (\shift cm,-2.35cm)
    {\(v_0v_{\j}\): yes\\\txt};
}
\end{tikzpicture}}
\newcommand{\figsplit}{%
\begin{tikzpicture}[scale=.67]
\draw[polyA] (0,-3) rectangle (6,4);
\draw[polyB] (-3,-2)--(8,3.5)--(8,6)--(-3,6)--cycle;
\node[smalllabel] at (5,-1.5) {\(P_i\)};
\node[smalllabel] at (-1.5,4.5) {\(P_h\)};
\foreach \x/\y in {0/-.5,6/2.5} {\fill[bad] (\x,\y) circle (3pt);}
\draw[very thick,good] (0,-.5)--(6,2.5);
\node[smalllabel] at (3,-3.6) {proper crossings split the lower edge};
\begin{scope}[xshift=11cm]
\draw[polyA] (0,0) rectangle (5,2);
\draw[polyB] (2,0) rectangle (7,3);
\node[smalllabel] at (1,1) {\(P_i\)};
\node[smalllabel] at (6,2) {\(P_h\)};
\draw[very thick,good] (2,0)--(5,0);
\fill[bad] (2,0) circle (3pt); \fill[bad] (5,0) circle (3pt);
\node[smalllabel] at (3.5,-.7) {shared-edge overlap: two endpoint events};
\end{scope}
\end{tikzpicture}}
\newcommand{\figlimits}{%
\begin{tikzpicture}[scale=.89]
\begin{scope}
\fill[pOne!8] (0,-4.5) rectangle (1.4,1.3);\draw[pOne,thick] (0,-4.5)--(0,1.3);
\draw[pTwo,thick] (-3,-1.5)--(1.4,.7);\draw[route,-Latex] (-3,-4)--(1,.5);
\node[smalllabel] at (.7,-3) {\(P_1\)};
\node[smalllabel] at (-2,.1) {\(\ell_2\)};
\fill[ink] (-3,-4) circle (1.8pt); \fill[ink] (1,.5) circle (1.8pt);
\node[smalllabel,below left] at (-3,-4) {\(s\)};
\node[smalllabel] at (-1,-5) {\(q_+(1/2)=(1,1/2)\)};
\end{scope}
\begin{scope}[xshift=6.5cm]
\fill[pOne!8] (0,-4.5) rectangle (1.4,1.3);\draw[pOne,thick] (0,-4.5)--(0,1.3);
\draw[pTwo,thick] (-3,-1.5)--(1.4,.7);\draw[route,-Latex] (-3,-4)--(0,-11/8)--(-1,-.5);
\node[smalllabel] at (.7,-3) {\(P_1\)};
\node[smalllabel] at (-2,.1) {\(\ell_2\)};
\fill[ink] (-3,-4) circle (1.8pt); \fill[ink] (-1,-.5) circle (1.8pt);
\node[smalllabel,below left] at (-3,-4) {\(s\)};
\fill[ink] (0,-11/8) circle (1.5pt);
\node[smalllabel] at (-1,-5) {\(q_-(1/2)=(-1,-1/2)\)};
\end{scope}
\end{tikzpicture}}
\newcommand{\figrays}{%
\begin{tikzpicture}[scale=1.15]
\coordinate (J) at (0,0);
\draw[polyA] (0,-3)--(0,2); \draw[polyB] (-3,-1.5)--(3,1.5);
\node[smalllabel] at (.3,1.7) {\(\ell_1\)};
\node[smalllabel] at (-2,-1.2) {\(\ell_2\)};
\draw[wrong,-Latex] (J)--(-1.8,2.4) node[above] {\(R_1u\)};
\draw[wrong,-Latex] (J)--(3,0) node[right] {\(R_2u\)};
\draw[route,-Latex] (J)--(.875,-3) node[below] {\(R_2R_1u\)};
\draw[route,-Latex] (J)--(2.6,0);
\fill[ink] (J) circle (2pt);\node[smalllabel,above right] at (J) {\(j\)};
\fill[good] (2,-1.5) circle (2pt);\node[smalllabel,right] at (2,-1.5) {\(t_A/2\)};
\fill[bad] (-2,-1.5) circle (2pt);\node[smalllabel,left] at (-2,-1.5) {\(t_B/2\)};
\end{tikzpicture}}
\newcommand{\figvirtual}{%
\begin{tikzpicture}[scale=1.15]
\draw[polyA] (0,-2)--(0,2);\draw[polyB] (-3,-1.5)--(3,1.5);
\node[smalllabel] at (.3,1.7) {\(\ell_1\)};
\node[smalllabel] at (-2,-1.2) {\(\ell_2\)};
\coordinate (S) at (-1.5,-2);\coordinate (V1) at (1.5,-2);\coordinate (V2) at (-.7,2.4);
\fill[ink] (S) circle (2pt);\fill[ink] (V1) circle (2pt);\fill[ink] (V2) circle (2pt);
\node[smalllabel,below] at (S) {\(s\)};\node[smalllabel,below] at (V1) {\(R_1s\)};
\node[smalllabel,above] at (V2) {\(R_2R_1s\)};
\draw[aux,-Latex] (S)--(V1);\draw[aux,-Latex] (V1)--(V2);
\fill[good] (-2,-1.5) circle (2pt);\node[smalllabel,left] at (-2,-1.5) {\(t_B/2\)};
\end{tikzpicture}}
\newcommand{\figpath}{%
\begin{tikzpicture}[scale=.95]
\begin{scope}
\clip (-5,-4.6) rectangle (5,1.4);
\draw[polyB] (-8,-4)--(8,4)--(8,9)--(-8,9)--cycle;
\draw[polyA] (0,-6) rectangle (6,6);
\node[smalllabel] at (4,-2) {\(P_1\)};
\node[smalllabel] at (-2,.5) {\(P_2\)};
\draw[route,-Latex] (-3,-4)--(0,-3)--(-3.6,-1.8)--(-4,-3);
\draw[wrong] (-3,-4)--(0,0)--(-4,-3);
\end{scope}
\foreach \x/\y/\lab in {-3/-4/s,0/-3/a,-3.6/-1.8/b,-4/-3/t_B,0/0/j}
  {\fill[ink] (\x,\y) circle (2pt);\node[smalllabel,above right] at (\x,\y) {\(\lab\)};}
\node[smalllabel] at (2.3,-4.1) {solid: returned route};
\end{tikzpicture}}
\newcommand{\figbackwards}{%
\begin{tikzpicture}[scale=1.14]
\draw[aux] (-3,0)--(4,0);
\draw[good,-Latex,line width=1.3pt] (0,0)--(3,0);
\draw[bad,-Latex,line width=1.3pt] (0,0)--(-2,0);
\fill[ink] (0,0) circle (2pt); \fill[bad] (-2,0) circle (2pt);
\node[smalllabel,above] at (0,0) {\(v\)};
\node[smalllabel,above] at (2,0) {forward ray \(r\)};
\node[smalllabel,above] at (-2,0) {\(q\) behind \(v\)};
\node[smalllabel,below] at (0,-.4) {\(r\times(q-v)=0\), but \(r\cdot(q-v)<0\)};
\end{tikzpicture}}
\newcommand{\figmembership}{%
\begin{tikzpicture}[scale=1.15]
\fill[pOne!14] (0,0)--(4,2)--(2,4)--cycle;
\draw[polyA,fill opacity=.25] (0,0)--(3,0)--(4,2)--(2,4)--(-1,2)--cycle;
\node[smalllabel] at (2.7,2.8) {\(P_i\)};
\draw[aux] (0,0)--(4,2); \draw[aux] (0,0)--(2,4);
\fill[good] (2,1.5) circle (2.5pt);
\foreach \x/\y/\n in {0/0/0,3/0/1,4/2/2,2/4/3,-1/2/4}
  {\fill[ink] (\x,\y) circle (1.5pt);\node[smalllabel,above right] at (\x,\y) {\(p_{\n}\)};}
\node[smalllabel,right] at (2,1.5) {\(q\)};
\node[smalllabel] at (2,-.55) {fan triangle \((p_0,p_2,p_3)\)};
\end{tikzpicture}}

% A sign decision is a coefficient comparison, not a numerically tiny offset.
\newcommand{\figsymbolic}{%
\begin{tikzpicture}[font=\small,
 stage/.style={draw=navy!35,fill=codebg,rounded corners=2pt,
                 text width=2.8cm,minimum height=1.35cm,align=center,inner sep=4pt}]
\node[stage] (a) at (0,0) {\(\underbrace{a\times(q-o)}_{0}\)\\constant};
\node[stage] (b) at (3.55,0) {\(\underbrace{a\times d}_{0}\)\\\(\varepsilon\) side};
\node[stage] (c) at (7.1,0) {\(\underbrace{a\times e_1}_{0}\)\\\(\varepsilon^2\) tie};
\node[stage,draw=good,fill=good!7] (d) at (10.65,0)
  {\(\underbrace{a\times e_2}_{+1}\)\\\(\varepsilon^3\) decides};
\foreach \u/\v in {a/b,b/c,c/d} \draw[quiet,-Latex] (\u.east)--(\v.west);
\node[font=\scriptsize,color=quiet,align=center] at (5.3,-1.25)
  {\(a=(1,0),\ o=(0,0),\ q=(2,0),\ d=(-1,0),\ e_1=(1,0),\ e_2=(0,1)\)};
\end{tikzpicture}}

% Two rows distinguish recursive query transforms from physical refolding.
\newcommand{\figrefoldtrace}{%
\begin{tikzpicture}[font=\small,
 tracebox/.style={draw=navy!35,fill=codebg,rounded corners=2pt,
              text width=3.65cm,minimum height=1.05cm,align=center,inner sep=4pt}]
\node[font=\scriptsize\bfseries,color=navy,anchor=east] at (-.35,0)
  {UNFOLD};
\node[tracebox] (q) at (2,0) {\(t_B=(-4,-3)\)\\level 2};
\node[tracebox] (qp) at (7,0) {\(q'=(-24/5,-7/5)\)\\level 1};
\node[tracebox] (qpp) at (12,0) {\(q''=(24/5,-7/5)\)\\level 0};
\draw[good,-Latex,line width=.9pt] (q.east)--node[above,font=\scriptsize]{\(R_{\ell_2}\)}(qp.west);
\draw[good,-Latex,line width=.9pt] (qp.east)--node[above,font=\scriptsize]{\(R_{\ell_1}\)}(qpp.west);
\node[font=\scriptsize\bfseries,color=navy,anchor=east] at (-.35,-2.15)
  {REFOLD};
\node[tracebox] (p2) at (2,-2.15) {\(s\to a\to b\to t_B\)\\\(b=(-18/5,-9/5)\)};
\node[tracebox] (p1) at (7,-2.15) {\(s\to a\to q'\)\\\(a=(0,-3)\)};
\node[tracebox] (p0) at (12,-2.15) {\(s\to q''\)\\base path};
\draw[good,-Latex,line width=.9pt] (p0.west)--node[above,font=\scriptsize]{\(\ell_1\)}(p1.east);
\draw[good,-Latex,line width=.9pt] (p1.west)--node[above,font=\scriptsize]{\(\ell_2\)}(p2.east);
\draw[quiet,-Latex] (qpp.south)--(p0.north);
\end{tikzpicture}}
